\subsection{Text assembly and merging}\label{subsec:doc-text-assembly}

Text assembly converts the cleaned layout into text. It consumes the output of two-pass extraction of Subsection~\ref{subsec:doc-two-pass}: a cleaned version of PDF from which headers, footers, sidebars, figures, tables, and invisible texts are removed. Therefore, what reaches the text assembly stage is expected to be body text and section heading elements, in reading order. The purpose of the text assembly stage is to turn these elements into hierarchical paragraphs that the knowledge extraction pipeline expects. 

\pinktodo{Hierarchical paragraph claim contradicts the sentence-only structure of knowledge extraction input, change that.}

\pinktodo{Remove "and so on" if that is a problem.}

Text assembly stage begins with a small filtering step to remove any survivors that passed the two-pass filter, but should not have: page numbers, page headers that could not be removed,journal metadata such as "Received", "Accepted", or publication date, and so on. After the filtering step, the remaining Docling elements are traversed in reading order. Section headers and their nesting levels are tracked during this traversal, so that each paragraph has a section path under which it appears (e.g. Methods > 2.1. Staining). Duplicate blocks found within the same section path are removed. 

The main operation in this stage is paragraph merging. In journal articles, it was observed that paragraphs are split across pages or columns because of page breaks, page headers or footers, or intervening figures and tables. The layout parser may return some of these fragments as separate blocks, even though they belong to the same paragraph. A context-aware stitcher rejoins these fragments: a block is considered to continue to the next narrative block, when it ends without clear sentence-ending closure. This includes trailing hyphens or commas, connectives such as "and" or "of", and mid-sentence abbreviations such as "et al." or "Fig.". In these cases, the two narrative blocks are merged. 

Citation markers such as "[1,2,3]", double-spaces, URLs are detected with regular expressions, and are removed before merging. This removal prevents these markers from interrupting the paragraph merging operation, as an intervening citation marker might prevent two blocks to be merged, which otherwise correctly would have been. 

The merging rule is purely syntactic: it only examines tokens and punctuation, not sentence meaning. Therefore, it can over-merge or not merge when it should. This is an accepted trade-off. The heuristic recovers common page-break and column-break cases without adding a second semantic parser. 

The result is a list of hierarchical rows, consisting of merged paragraph text, section path, and the position of the text within its document section. These rows are the textual units, over which the knowledge-extraction pipeline operates. Assigning each row to its position within the section, linking the figures and tables it cites, and writing into the database are explained in Subsection~\ref{subsec:doc-provenance-preserving-storage}. 


